%!TEX root = ../main.tex

\chapter{Preliminaries \label{pre}}
This chapter covers important theoretical aspects relevant for this thesis.
Section \ref{sec:ocp} formulates the \ac{OCP} of drone racing and presents a theoretical solution based on the application of \ac{PMP}.
The optimal solution derived from \ac{PMP} serves primarily as an upper theoretical bound, as many of its underlying assumptions do not hold in real-world scenarios.
Practical implementations must account for model uncertainties, disturbances, and system constraints that deviate from the idealized formulation.

In this context, Sections \ref{sec:classical} and \ref{sec:learned} review the state-of-the-art control approaches used to achieve high-performance drone racing, covering both classical control techniques and modern learning-based methods.
\section{Optimal Control Problem} \label{sec:ocp} 
We refer the reader to~\cite{lewis_optimal_2012} 


\subsection{General Discrete-Time Optimal Control Problem}
In general, a nonlinear discrete-time \ac{OCP} can be formulated as the optimization of a running cost (or loss) function $L_{k}()$, plus a terminal cost function $\phi()$, subject to the system's dynamics $\mathbf{f}_k()$, and the set of permitted states $\mathcal{X}_k$, and controls $\mathcal{U}_k$. 


\begin{equation}
\begin{aligned}
% & \underset{\boldsymbol{u}_{i}, \dots, \boldsymbol{u}_{N-1}, \boldsymbol{x}_{i}, \dots, \boldsymbol{x}_{N}}{\text{minimize}} & & \left( \phi(\boldsymbol{x}_{N}, N) + \sum_{k=i}^{N-1} L_{k}(\boldsymbol{x}_{k}, \boldsymbol{u}_{k}) \right) \\
& \underset{\mathbf{X},\ \mathbf{U}}{\text{minimize}} & & \left( \phi(\boldsymbol{x}_{N}, N) + \sum_{k=0}^{N-1} L_{k}(\boldsymbol{x}_{k}, \boldsymbol{u}_{k}) \right) \\
& \text{subject to} & & \boldsymbol{x}_{k+1} = \mathbf{f}_{k}(\boldsymbol{x}_{k}, \boldsymbol{u}_{k}), \quad k = 0, \dots, N-1, \\
& & & \boldsymbol{u}_{k} \in \mathcal{U}_{k}, \quad k = 0, \dots, N-1, \\
& & & \boldsymbol{x}_{k} \in \mathcal{X}_{k}, \quad k = 0, \dots, N,
\end{aligned}
\end{equation}


\noindent Where
\begin{itemize}
\item[] $N$ is the final discrete time,
\item[] $\boldsymbol{x}_k \in \mathbb{R}^n$ is the state at the discrete time $k \in \mathbb{Z}$,
\item[] $\boldsymbol{u}_k \in \mathbb{R}^m$ is the control at the discrete time $k \in \mathbb{Z}$,
\item[] $\mathbf{X} = [\boldsymbol{x}_{0}, \dots, \boldsymbol{x}_{N}]$ is the trajectory of states,
\item[] $\mathbf{U} = [\boldsymbol{u}_{0}, \dots, \boldsymbol{u}_{N-1}]$ is the trajectory of control actions.

\end{itemize}

\subsection{Minimum-Time Optimal Control Problem} \label{subsec:minOCP}
The last formulation of an \ac{OCP} can be further specified for the case of drone racing in several manners.
Since the main objective is to complete a lap through the racing track as fast as possible, this problem can be written as minimum-time problem.


\begin{equation}
\begin{aligned}
% & \underset{\boldsymbol{u}_{i}, \dots, \boldsymbol{u}_{N-1}, \boldsymbol{x}_{i}, \dots, \boldsymbol{x}_{N}}{\text{minimize}} & & \left( \phi(\boldsymbol{x}_{N}, N) + \sum_{k=i}^{N-1} L_{k}(\boldsymbol{x}_{k}, \boldsymbol{u}_{k}) \right) \\
 & \underset{\mathbf{U},\ N}{\text{minimize}} & & \qquad \qquad \sum_{k=0}^{N-1} 1 \\
 & \text{subject to} & & \boldsymbol{x}_{k+1} = \mathbf{f}_{k}(\boldsymbol{x}_{k}, \boldsymbol{u}_{k}), \quad k = 0, \dots, N-1, & \\
& & & \boldsymbol{u}_{k} \in \mathcal{U}_{k}, \quad k = 0, \dots, N-1, \\
& & & \boldsymbol{x}_{k} \in \mathcal{X}_{k}, \quad k = 0, \dots, N,
\end{aligned}
\end{equation}

The racing track restrictions, i.e. passing through a sequence of gates while avoiding
obstacles, can be introduced as hard constraints under the set of permitted states
$\mathcal{X}_k$ as

\begin{equation}
    \mathcal{X}_k = \left\{ \boldsymbol{x}_k \in \mathbb{R}^n \;\middle|\;
        \boldsymbol{p}_k \in \mathcal{G}_{g(k)}, \quad
        \boldsymbol{p}_k \notin \mathcal{O}_j, \quad
        \forall\, j \in \{1, \dots, M\}
    \right\},
\end{equation}

\noindent where
\begin{itemize}
    \item[] $\boldsymbol{p}_k \in \mathbb{R}^3$ is the position of the drone at time step $k$,
    \item[] $\mathcal{G}_{g(k)}$ is the feasible region defined by gate $g$ at time step $k$,
            with $g(k) \in \{1, \dots, G\}$ denoting the index of the gate to be passed
            at $k$,
    \item[] $G$ is the total number of gates in the racing track,
    \item[] $\mathcal{O}_j$ is the region occupied by the $j$-th obstacle,
    \item[] $M$ is the total number of obstacles.
\end{itemize}

\subsection{Simplified Minimum-Time Optimal Control Problem}~\label{sec:simpmin-OCP}
The original problem formulation (\autoref{subsec:minOCP}) is inherently complex, as it requires simultaneously solving for an optimal trajectory that clears all gates and avoids obstacles while determining the optimal control sequence achievable within the vehicle's dynamics and actuator constraints.
To make the problem more tractable, this thesis assumes that a reference trajectory or path is provided a priori.
Consequently, the \acp{OCP} addressed hereafter focus on trajectory tracking or path progress, formulated as follows:

\subsubsection{Trajectory Tracking}
In the trajectory tracking \ac{OCP}, the objective is to find the optimal control sequence $\mathbf{U}$ that minimizes the error between a given reference state $\boldsymbol{x}_k^{\text{ref}}$ and the current system state $\boldsymbol{x}_k$, subject to dynamics, actuator, and state constraints.
While this simplifies the optimization significantly it comes at the cost of reduced robustness.
By assuming the trajectory is fixed, the system cannot re-optimize the path to accommodate state constraints or environmental changes, effectively limiting the drone's agility compared to a unified spatial-temporal optimization.

\begin{equation}
\begin{aligned}
 & \underset{\mathbf{U}}{\text{minimize}} & & \sum_{k=0}^{N-1} \left\| \boldsymbol{x}_k - 
        \boldsymbol{x}_k^{\text{ref}} \right\|_{\mathbf{Q}}^{2} \\%+         \left\| \boldsymbol{u}_k \right\|_{\mathbf{R}}^{2} \\
 & \text{subject to} & & \boldsymbol{x}_{k+1} = \mathbf{f}_{k}(\boldsymbol{x}_{k}, 
        \boldsymbol{u}_{k}), \quad k = 0, \dots, N-1, \\
& & & \boldsymbol{u}_{k} \in \mathcal{U}_{k}, \quad k = 0, \dots, N-1, \\
& & & \boldsymbol{x}_{k} \in \mathcal{X}_{k}, \quad k = 0, \dots, N,
\end{aligned}
\end{equation}


% \subsubsection{Path Progress}
% An alternative to reduce the complexity without sacrifiyng much on robsutness is to optimize the progress along a given path. 
% Unlike trajectory tracking, which imposes a strict temporal constraint for every spatial coordinate, the path progress \ac{OCP} treats the speed along the path as a decision variable. 
% This allows the controller to dynamically adjust the vehicle's velocity based on current tracking errors or actuator saturation, ensuring the drone stays on the geometric path even when pushed to its physical limits.

% \section{Pontryagin's Maximum Principle}
% \todo{I'll probably remove this section...}

%  maybe solve it using a pmm and Pontryagin's Maximum Principle

% commment on: Since it is impossible to perfectly model the state transition equation, and the state estimation is not perfect, 
% the optimal solution breaks under real world conditions.

\section{Classical Control approaches}\label{sec:classical}

% \todo{Don't dive deep in these, just comment relevant stuff}
\subsection{Differential Flatness Based Control}
\subsection{Model Predictive Control}

\section{Learning-Based approaches}\label{sec:learned}
Learning-based approaches can be understood as optimization methods in which the optimization process is done by intereacting with the environment~\cite{huang_model-based_2020}.
There are two main types of \ac{RL}, model-based and model-free, each have with their own advantages and drawbacks.
Next we give an overview model-free and model-based \ac{RL} and the most relevant algorithms in the literature of each.
% \subsection{Neural Networks?}
\subsection{Model-Free Reinforcement Learning}
In model-free \ac{RL}, the neither the environment nor the dynamics of the agent are taking into account for the optimization process.
Instead, the policy is optimized by gathering data from interacting with the environment. 
The two main methods to do this, are Q-Learning and Policy Optimization.
% \subsubsection{Q-Learning}
% In Q-Learning the idea is to map the values of the possible actions given the current state and store those values into a table.
% Druing training, Bellman's equation is used to update the values of the Q-table, and when training has converged an optimal policy $\pi^*$ has been found since dor each state the table indicates the best action~\cite{Sutton1998}.
% \begin{equation}
%     \pi^*(s) = \arg\max_a Q^*(s,a)
% \end{equation}
% Q-Learning is a formulation of \ac{DP}, for smiple scenarios such us finindg the optimal policy to exit a maze it works, howerver for more complex environments it is computationally coslty and and unfesible approach.
% To relax the computational cost, \ac{DQN} was proposed in~\cite{mnih_playing_nodate}, where instead of storing the values in a Q table, a neural network is optimized to approximate the values.


% In model-free \ac{RL}, neither the environment model nor the agent dynamics are explicitly considered during the optimization process.
% Instead, the policy is optimized by collecting data through interactions with the environment.
% The two main approaches for doing this are Q-learning and policy optimization.

\subsubsection{Q-Learning}

In Q-learning, the objective is to estimate the value of taking a given action at a particular state and store these values in a Q-table.
During training, Bellman’s equation is used to iteratively update the entries of the table.
Once training converges, an optimal policy $\pi^*$ can be obtained, since the Q-table provides the action with the highest expected return for each state~\cite{Sutton1998}.

\begin{equation}
    \pi^*(s) = \arg\max_a Q^*(s,a)
\end{equation}

Q-learning can be interpreted as a formulation of \ac{DP}.
For simple problems, such as finding the optimal strategy to escape a maze, this approach performs well.
However, for high-dimensional or complex environments, storing and updating the Q-table becomes computationally infeasible.

To overcome this limitation, \ac{DQN} was introduced in~\cite{mnih_playing_nodate}.
Instead of explicitly storing the Q-values in a table, a neural network is trained to approximate the action-value function.


\subsubsection{Policy Optimization}
In policy optimization methods, instead of learning the action-value function and deriving an optimal policy by selecting the best action at a given state, the policy is learned directly through the parameters of a neural network $\theta$~\cite{Sutton1998}.
These parameters are optimized with respect to an objective function $J(\theta)$ that encodes the desired behavior of the policy.

\begin{equation}\label{opt_theta}
    \theta^* = \arg \max_{\theta} \mathbb{E}_{\tau \sim \pi_\theta} \left[ R(\tau) \right]
\end{equation}

In this framework, policy optimization methods aim to find the optimal parameters $\theta^*$ such that the resulting policy $\pi_{\theta}(a|s)$ selects actions $a \in \mathcal{A}$ for states $s \in \mathcal{S}$ that maximize the expected return over trajectories $\tau = \{s_0, a_0, s_1, a_1, \dots\}$.


\subsection{Model-Based Reinforcement Learning}

\section{Differentiable Simulation}